% problem-type: truyen-song-3d-doi-xung-cau
% order: 40
% Bo cuc: De vi du -> Cong thuc -> De + Bai giai

\section{Truyền sóng 3D có nguồn --- đối xứng cầu ($v=ru$)}

% ---------------------------------------------------------------
\subsection*{Đề bài ví dụ}
% ---------------------------------------------------------------
Giải $u_{tt}=\Delta u$ trên $\mathbb{R}^3$, $u(\cdot,0)=0$, $u_t(\cdot,0)=\psi$ với
$\psi$ \textbf{đối xứng cầu}: $\psi=\chi_{B_1(0)}$ (bằng $1$ khi $r=|x|<1$). \textbf{Tính $u(0,t)$ tại tâm.}

\emph{\small(Ví dụ tự soạn minh hoạ riêng kỹ thuật đổi biến cầu $v=ru$; kỹ thuật này dùng cho phần ``đối xứng cầu'' của các đề sóng 3D có nguồn: MAT3365 HK2 2023--24 Câu 3, MAT3365 HK1 2025--26 Câu 3.)}

% ---------------------------------------------------------------
\subsection*{Nhận ra dạng này khi nào?}
% ---------------------------------------------------------------
\begin{itemize}
  \item Sóng 3D, dữ liệu/nguồn \textbf{chỉ phụ thuộc $r=|x|$} (đối xứng cầu).
  \item Khi đó dùng đổi biến $v=ru$ để hạ về sóng 1D (gọn hơn Kirchhoff).
\end{itemize}
\begin{meobox}
Dấu hiệu: $\psi,f$ viết theo $|x|$ (hoặc $\chi_{B_\rho}$). Tâm $r=0$: $u(0,t)=v_r(0,t)$.
\end{meobox}
\begin{luuybox}
Đề thực thường \textbf{TRỘN}: phần đối xứng cầu (mục này) $+$ phần đối xứng phẳng (phụ thuộc 1 biến).
Chồng chất tách rồi xử lý từng phần; phần phẳng xem mục ``đối xứng phẳng (hạ 1 biến)''.
\end{luuybox}

% ---------------------------------------------------------------
\subsection*{Công thức \& phương pháp}
% ---------------------------------------------------------------
\subsubsection*{Đổi biến cầu $v=ru$}
% technique: doi-bien-cau-song
\textbf{Đổi biến cầu $v=r\,u$ cho sóng đối xứng cầu.}
\emph{$u=u(r,t)$ là nghiệm cần tìm} (biên độ tại khoảng cách $r$ tới tâm); \emph{$v=r\,u$ là hàm phụ} đặt ra để tính cho dễ.

\textbf{Vì sao đặt $v=ru$:} Laplacian của hàm chỉ phụ thuộc $r$ có số hạng phiền $\tfrac{2}{r}u_r$:
\[
  \Delta u = u_{rr}+\tfrac{2}{r}u_r .
\]
Nhân $u$ với $r$: $v_r=u+ru_r$, $v_{rr}=2u_r+ru_{rr}=r\big(u_{rr}+\tfrac{2}{r}u_r\big)=r\,\Delta u$. Nhân phương trình sóng với $r$:
\[
  r\,u_{tt}=c^2 r\,\Delta u+rf \;\Longrightarrow\; v_{tt}=c^2 v_{rr}+rf .
\]
Số hạng $\tfrac{2}{r}u_r$ \textbf{biến mất} $\Rightarrow$ $v$ thỏa \textbf{sóng 1D phẳng} trên $r>0$ (có nguồn $rf$ nếu có nguồn).

\textbf{Giải:} dùng d'Alembert cho $v$, rồi lấy lại $u=v/r$. Biên $v(0,t)=ru\big|_{r=0}=0$ (vì $u$ hữu hạn ở tâm) là \textbf{Dirichlet} $\Rightarrow$ thác triển dữ liệu/nguồn \textbf{lẻ} qua $r=0$ (xem ``Thác triển chẵn/lẻ'').
Dùng khi dữ liệu/nguồn đối xứng cầu (vd $\psi,f$ chỉ phụ thuộc $r$).

\emph{(Có nguồn $f(r,t)$: $v=ru$ cho $v_{tt}=c^2v_{rr}+rf$, giải bằng d'Alembert 1D có nguồn.)}

% ---------------------------------------------------------------
\subsection*{Đề bài \& bài giải}
% ---------------------------------------------------------------
\paragraph{Nhắc lại đề.} $u_{tt}=\Delta u$ trên $\mathbb{R}^3$, $u(\cdot,0)=0$, $u_t(\cdot,0)=\psi=\chi_{B_1(0)}$
(đối xứng cầu). Tính $u(0,t)$ tại tâm.

\paragraph{Bước 1 --- đổi biến.} $c=1$, đặt $v=ru$. Vì $u(\cdot,0)=0$, $u_t(\cdot,0)=\psi(r)$:
\[
  v_{tt}=v_{rr}\ (r>0),\quad v(r,0)=0,\quad v_t(r,0)=r\,\psi(r)=r\,\chi_{[0,1]}(r),\quad v(0,t)=0.
\]

\paragraph{Bước 2 --- d'Alembert (thác triển lẻ qua $r=0$).} Với $h(r)=r\chi_{[0,1]}(r)$, thác triển lẻ
$h_{\mathrm{lẻ}}(s)=s\,\chi_{[-1,1]}(s)$:
\[
  v(r,t)=\tfrac12\int_{r-t}^{r+t} h_{\mathrm{lẻ}}(s)\,ds,\qquad
  u(r,t)=\frac{v(r,t)}{r}.
\]

\paragraph{Bước 3 --- tại tâm $r=0$ (khử dạng $0/0$).} $u=v/r$ tại $r=0$ là $0/0$ vì
$v(0,t)=\tfrac12\int_{-t}^{t}h_{\mathrm{lẻ}}=0$ ($h_{\mathrm{lẻ}}$ lẻ). L'Hôpital theo $r$:
$u(0,t)=\lim_{r\to0}\tfrac{v(r,t)}{r}=v_r(0,t)$.

\paragraph{Bước 3a --- đạo hàm theo cận.} $v_r(r,t)=\tfrac12\big[h_{\mathrm{lẻ}}(r+t)-h_{\mathrm{lẻ}}(r-t)\big]$.
Tại $r=0$, dùng $h_{\mathrm{lẻ}}(-t)=-h_{\mathrm{lẻ}}(t)$:
\[
  u(0,t)=\tfrac12\big[h_{\mathrm{lẻ}}(t)-h_{\mathrm{lẻ}}(-t)\big]=h_{\mathrm{lẻ}}(t).
\]

\paragraph{Bước 3b --- thay $h_{\mathrm{lẻ}}(t)=t\,\chi_{[-1,1]}(t)$.}\leavevmode

\begin{nhobox}
\[
  u(0,t)=\begin{cases} t, & 0<t<1,\\ 0, & t>1.\end{cases}
\]
Tín hiệu qua tâm trong $(0,1)$ rồi tắt --- khớp Huygens mạnh (so sánh: dạng bài Kirchhoff cho cùng kết quả).
\end{nhobox}